% Лекц 16: семинар ба гэрийн даалгавар. Бодлогууд: Г.Н. Берман, "Сборник задач по курсу
% математического анализа" (22-р хэвлэл). [Б.N] нь номын бодлогын дугаар.

\ProblemSet{Семинар 16}
\begin{frame}[allowframebreaks]{Семинар 16}
\pgroup{Хязгааргүй хилтэй өргөтгөсөн интегралыг бод (эсвэл сарнихыг тогтоо).}
\begin{plist}
\prob{2366}{$\int_1^{+\infty}\frac{dx}{x^4}$}
\prob{2367}{$\int_1^{+\infty}\frac{dx}{\sqrt x}$}
\prob{2368}{$\int_0^{+\infty}e^{-ax}dx$ \ ($a>0$)}
\prob{2369}{$\int_{-\infty}^{+\infty}\frac{2x\,dx}{x^2+1}$}
\prob{2370}{$\int_{-\infty}^{+\infty}\frac{dx}{x^2+2x+2}$}
\prob{2376}{$\int_0^{+\infty}xe^{-x^2}dx$}
\prob{2380}{$\int_0^{+\infty}e^{-x}\sin x\,dx$}
\end{plist}
\pgroup{Нийлэлтийг шинжил.}
\begin{plist}
\prob{2386}{$\int_0^{+\infty}\frac{x\,dx}{x^3+1}$}
\prob{2387}{$\int_1^{+\infty}\frac{x^3+1}{x^4}\,dx$}
\prob{2412}{$\int_0^1\frac{\sqrt x\,dx}{\sqrt{1-x^4}}$}
\end{plist}
\pgroup{Үл зааглагдах функцийн интегралыг бод (эсвэл сарнихыг тогтоо).}
\begin{plist}
\prob{2394}{$\int_0^1\frac{dx}{\sqrt{1-x^2}}$}
\prob{2395}{$\int_0^2\frac{dx}{x^2-4x+3}$}
\prob{2397}{$\int_0^1x\ln x\,dx$}
\prob{2399}{$\int_1^2\frac{dx}{x\ln x}$}
\end{plist}
\pgroup{Талбай ба эзлэхүүн.}
\begin{plist}
\prob{1600}{$y=x^2-6x+10$ ба $y=6x-x^2$ параболуудаар хүрээлэгдсэн дүрсийн талбайг ол.}
\prob{2460}{$y=x^2$ ба $y=\frac{x^3}{3}$ шугамуудаар хүрээлэгдсэн дүрсийн талбайг ол.}
\prob{2478}{$y=e^x$, $y=e^{-x}$ ба $x=1$ шулуунаар хүрээлэгдсэн дүрсийн талбайг ол.}
\prob{2559}{$y=xe^x$, $x=1$, $y=0$ шугамуудаар хүрээлэгдсэн муруй шугаман трапецийг $Ox$ тэнхлэгийг тойруулан эргүүлэхэд үүсэх биетийн эзлэхүүнийг ол.}
\prob{2561}{$y=x^2$ ба $y^2=x$ параболуудаар хүрээлэгдсэн дүрсийг $Ox$ тэнхлэгийг тойруулан эргүүлэхэд үүсэх биетийн эзлэхүүнийг ол.}
\prob{2564}{$y=2x-x^2$ парабол ба $Ox$ тэнхлэгээр хүрээлэгдсэн дүрсийг $Oy$ тэнхлэгийг тойруулан эргүүлэхэд үүсэх биетийн эзлэхүүнийг ол.}
\end{plist}
\end{frame}

\ProblemSet{Гэрийн даалгавар 16}
\begin{frame}[allowframebreaks]{Гэрийн даалгавар 16}
\pgroup{Хязгааргүй хилтэй өргөтгөсөн интегралыг бод (эсвэл сарнихыг тогтоо).}
\begin{plist}
\prob{2371}{$\int_2^{+\infty}\frac{\ln x}{x}\,dx$}
\prob{2372}{$\int_1^{+\infty}\frac{dx}{x^2(x+1)}$}
\prob{2373}{$\int_0^{+\infty}\frac{x\,dx}{(1+x)^3}$}
\prob{2377}{$\int_0^{+\infty}x^3e^{-x^2}dx$}
\prob{2378}{$\int_0^{+\infty}x\sin x\,dx$}
\prob{2381}{$\int_0^{+\infty}e^{-ax}\cos bx\,dx$ \ ($a>0$)}
\prob{2382}{$\int_1^{+\infty}\frac{\arctg x}{x^2}\,dx$}
\prob{2383}{$\int_0^{+\infty}\frac{dx}{1+x^3}$}
\prob{2385}{$\int_1^{+\infty}\frac{\sqrt x\,dx}{(1+x)^2}$}
\end{plist}
\pgroup{Нийлэлтийг шинжил.}
\begin{plist}
\prob{2389}{$\int_1^{+\infty}\frac{\ln(x^2+1)}{x}\,dx$}
\prob{2390}{$\int_0^{+\infty}\sqrt x\,e^{-x}dx$}
\prob{2392}{$\int_{e^2}^{+\infty}\frac{dx}{x\ln\ln x}$}
\prob{2393}{$\int_e^{+\infty}\frac{dx}{x(\ln x)^{3/2}}$}
\prob{2413}{$\int_0^1\frac{x^2\,dx}{\sqrt[3]{(1-x^2)^5}}$}
\prob{2414}{$\int_0^1\frac{dx}{e^{\sqrt x}-1}$}
\end{plist}
\pgroup{Үл зааглагдах функцийн интегралыг бод (эсвэл сарнихыг тогтоо).}
\begin{plist}
\prob{2398}{$\int_0^{1/e}\frac{dx}{x\ln^2x}$}
\prob{2400}{$\int_1^e\frac{dx}{x\sqrt{\ln x}}$}
\prob{2402}{$\int_a^b\frac{x\,dx}{\sqrt{(x-a)(b-x)}}$ \ ($a<b$)}
\prob{2405}{$\int_{-1}^1\frac{dx}{(2-x)\sqrt{1-x^2}}$}
\prob{2406}{$\int_{-1}^1\frac{3x^2+2}{\sqrt[3]{x^2}}\,dx$}
\prob{2411}{$\int_0^1\frac{e^{1/x}}{x^3}\,dx$}
\end{plist}
\pgroup{Талбай.}
\begin{plist}
\prob{1601}{$y=x^2-2x+2$ парабол, түүний $(3,5)$ цэг дээрх шүргэгч, $Oy$ ба $Ox$ тэнхлэгүүдээр хүрээлэгдсэн дүрсийн талбайг ол.}
\prob{2459}{$y^2+8x=16$ ба $y^2-24x=48$ параболуудаар хүрээлэгдсэн дүрсийн талбайг ол.}
\prob{2468}{$y=x(x-1)^2$ шугам ба абсциссын тэнхлэгээр хүрээлэгдсэн дүрсийн талбайг ол.}
\prob{2483}{$y=\ln x$ ба $y=\ln^2x$ шугамуудаар хүрээлэгдсэн дүрсийн талбайг ол.}
\prob{2485}{Абсциссын тэнхлэг, $y=\sin x$ ба $y=\cos x$ шугамуудаар хүрээлэгдсэн муруй шугаман гурвалжнуудын нэгийн талбайг ол.}
\end{plist}
\pgroup{Эзлэхүүн.}
\begin{plist}
\prob{2555}{$y^2=4x$ параболыг тэнхлэгийг нь тойруулан эргүүлэхэд үүсэх гадаргуу ба оройгоос 1 нэгжийн зайд, тэнхлэгт перпендикуляр хавтгайгаар хүрээлэгдсэн биетийн эзлэхүүнийг ол.}
\prob{2556}{Их тэнхлэг нь $2a$, бага тэнхлэг нь $2b$ эллипсийг 1) их тэнхлэгийг нь, 2) бага тэнхлэгийг нь тойруулан эргүүлэв. Үүсэх эллипсоидуудын эзлэхүүнийг ол. Тухайн тохиолдолд бөмбөрцгийн эзлэхүүнийг гарга.}
\prob{2563}{$y=\arcsin x$ шугамаар хүрээлэгдсэн, $[0,1]$ суурьтай муруй шугаман трапецийг $Ox$ тойруулан эргүүлэхэд үүсэх биетийн эзлэхүүнийг ол.}
\prob{2565}{$y=\sin x$ синусоидын хагас үеийн нумаар хүрээлэгдсэн муруй шугаман трапецийг $Oy$ тэнхлэгийг тойруулан эргүүлэхэд үүсэх биетийн эзлэхүүнийг ол.}
\end{plist}
\end{frame}
